% 自测题9：计算题9.3.1—9.3.5，证明题9.4.1；
% 源书印刷204—205，源PDF物理212—213页。
\subsection{自测题9：双缝、薄膜和牛顿环计算题}

\sourcetest{9.3.1}{204}{212}
\problemtext 杨氏双缝间距0.20mm、屏距1.0m，
入射400—760nm白光。
屏上离零级亮纹20mm的位置，哪些波长在该点最大增强？
\workings
几何光程差
\[
\Delta=\frac{dx}{D}
=\frac{(0.20\times10^{-3})(20\times10^{-3})}{1.0}
=4.0\times10^{-6}\,\mathrm m=4000\,\mathrm{nm}.
\]
第$m$级亮纹需$m\lambda=\Delta$，
于是$\lambda=4000/m$nm。
要求$400\le\lambda\le760$nm，
即$4000/760\le m\le4000/400$，
整数$m=6,7,8,9,10$。因此
\[
\boxed{\lambda\simeq666.7,\ 571.4,\ 500.0,\ 444.4,\ 400.0\ \mathrm{nm}}.
\]
\crosscheck 400nm波长在题给可见光端点，故应计入；
若只考虑严格开区间，需明确排除边界才可少计一个。

\sourcetest{9.3.2}{205}{213}
\problemtext 两缝照射波长546.1nm的平面单色光，
屏距$D=2.00$m；
零级亮纹两侧的第5级亮纹相距12.0mm。
求(1)缝距；(2)从零级向一侧数至第20条亮纹的距离；
(3)若平面光斜入射，条纹间隔如何变化？
\workings
设相邻亮纹间距$\beta=\lambda D/d$。
两侧级次为$m=\pm5$，所以$10\beta=12.0$mm，
\[
\beta=1.20\,\mathrm{mm}.
\]
由此
\[
\boxed{d=\frac{\lambda D}{\beta}
=\frac{(546.1\times10^{-9})(2.00)}
{1.20\times10^{-3}}
\simeq9.102\times10^{-4}\,\mathrm m=0.910\,\mathrm{mm}}.
\]
从零级起第20级的坐标大小为
\[
\boxed{x_{20}=20\beta=24.0\,\mathrm{mm}}.
\]
平面光斜入射添加固定的初始两缝相差$d\sin i$，
不改变相邻级次的几何路程差递增$\lambda$，
故在小角度理想条件下
\[
\boxed{\beta'=\beta=1.20\,\mathrm{mm}}.
\]
\crosscheck 屏上$\pm5$亮纹之间有\emph{10}个条纹间隔，
不是5个；斜入射移动整个图样的零点，而非必改变间距。

\sourcetest{9.3.3}{205}{213}
\problemtext 折射率$n_f=1.35$的透明膜镀于$n_s=1.50$玻璃上，
垂直入射光在600nm处反射相消、700nm处反射相长，
且600—700nm之间不存在其他反射极大或极小。
求薄膜厚度。
\workings
空气$\to$膜和膜$\to$玻璃两次反射都产生$\pi$相变，
相对半波损失为零。
记膜内双程光程$L=2n_ft$；
反射相长要求$L=m\lambda$，相消要求
$L=(k+\tfrac12)\lambda$。
600nm相消与700nm相长是\emph{相邻}极值，
故两者的等效级数相差半级：
\[
 L=(m+\tfrac12)(600\,\mathrm{nm})
 =m(700\,\mathrm{nm}).
\]
解$600m+300=700m$得$m=3$，
\[
\boxed{L=2100\,\mathrm{nm}},\qquad
\boxed{t=\frac{2100}{2(1.35)}\,\mathrm{nm}
\simeq777.8\,\mathrm{nm}}.
\]
\crosscheck 700nm对应相长第3级；
600nm对应相消的3.5级，
二者中间无其他半整数级极值。若误判一束反射
才有相变，答案可能错半个级次。

\sourcetest{9.3.4}{205}{213}
\problemtext 空气楔尖两板一端相接，单色光500nm垂直照射；
楔角$\alpha=2\times10^{-4}$rad。
向楔尖充入折射率$n=1.40$的液体，
求从棱边数起第5条\emph{明纹}移动距离和方向。
\workings
对玻璃板间低于玻璃折射率的介质膜，
两主要反射光有一次相对半波损失；
从接触处暗纹起第一明纹对应
$2nt=(1/2)\lambda$。
从棱边数起第5条明纹则有
$2nt=(5-\tfrac12)\lambda=4.5\lambda$。
位置为
$x_5(n)=4.5\lambda/(2n\alpha)$。
介质由空气变为$n=1.40$后
\[
\boxed{|\Delta x|
=\frac{4.5\lambda}{2\alpha}\left(1-\frac1n\right)
=\frac{4.5(500\,\mathrm{nm})}{4\times10^{-4}}
 \left(1-\frac1{1.4}\right)
\simeq1.607\,\mathrm{mm}}.
\]
移动方向为\textbf{向劈尖棱边（薄端）}。
\crosscheck 反射等厚纹接触处为暗，
若把第5条明纹当成“5个完整波长”将有半级误差。
液体只改变光学厚度，不改变题给楔角的几何条件。

\sourcetest{9.3.5}{205}{213}
\problemtext 原书图9.3.1中的牛顿环，平凸透镜曲率半径不变，
玻璃间原为空气$n=1.00$，
后来改为水$n'=1.33$。
求第$k$级暗环半径的相对缩小量$(r_k-r'_k)/r_k$。
\workings
同一环级反射暗纹半径
$r_k^2=kR\lambda_0/n$，故
\[
 \frac{r'_k}{r_k}=\sqrt{\frac{n}{n'}}
 =\frac1{\sqrt{1.33}}.
\]
于是
\[
\boxed{\frac{r_k-r'_k}{r_k}
=1-\frac1{\sqrt{1.33}}\simeq0.1329=13.29\%}.
\]
\crosscheck 第$k$级数和曲率半径同时消去，
因此相对缩小率不依赖$k$或$R$。

\subsection{自测题9：干涉场的振幅证明}

\sourcetest{9.4.1}{205}{213}
\problemtext 杨氏双缝两光束在屏幕P点振动为
$E_1=E_0\sin\omega t$及
$E_2=E_0\sin(\omega t+\Phi)$，
证明合振幅可写成
$E_P=E_m\cos[(\pi d/\lambda)\sin\theta]$，
其中$E_m$是最大合振幅。
\workings
利用和差化积恒等式
\[
\begin{aligned}
E_P&=E_1+E_2\\
&=2E_0\cos\frac\Phi2\,
 \sin\left(\omega t+\frac\Phi2\right).
\end{aligned}
\]
因此有向振幅系数为$2E_0\cos(\Phi/2)$，
最大幅值$E_m=2E_0$。
双缝缝距$d$，观测方向$\theta$时几何光程差
$\delta=d\sin\theta$，若双缝初相相同，
\[
\Phi=\frac{2\pi\delta}{\lambda}
=\frac{2\pi d}{\lambda}\sin\theta.
\]
于是
\[
\boxed{E_P=E_m\cos\left(\frac{\pi d}{\lambda}\sin\theta\right)}
\]
表示\emph{有向合振幅的系数}；
实际非负幅值应写
$|E_m\cos[(\pi d/\lambda)\sin\theta]|$。
\crosscheck 合成振动还包含时间相位
$\sin(\omega t+\Phi/2)$，
不能把有向振幅系数误认为完整瞬时电场。
